\section{Moving a Field That Is Already Answering}
\label{sec:ch12-moving-a-field-that-answers}
\index{override@\code{@override}|(}%

Everything so far assumed a field can be picked up in one service and put down
in another between two deployments. Sometimes it can. When it cannot, because
the field is answering traffic and the two services will be running side by side
for a while, \code{@override} is the directive for it, and it is the only one in
Apollo Federation whose whole purpose is a migration rather than a shape.

The declaration goes on the subgraph that is \emph{taking} the field:

\begin{graphqlsdl}[fontsize=\footnotesize]
type Product @key(fields: "id") {
  id: ID!
  availableQuantity: Int! @override(from: "catalog")
}
\end{graphqlsdl}

The subgraph it is taken from changes nothing at all. It keeps declaring the
field, keeps resolving it if asked, and stops being asked.

\index{override@\code{@override}!cases derived from the real schemas}%
Every composer message in this section is asserted by
\code{scripts/override-cases.mjs}, which both verification scripts run. Nine
cases, each one the six committed schemas with literal edits that have to match
exactly once, composed with wgc 0.129.7. One thing about them needs saying
plainly rather than being discovered later. At tag \code{ch12} no two subgraphs
declare the same field, because the extraction is finished, so a case that needs
a \enquote{before} has to invent one: five of the nine add
\code{availableQuantity} back to \code{catalog}, which stands in for the
monolith, and take it away again from \code{inventory}. The other four need only
one declaration to make their point and leave Catalog alone. The edit is
fictional. Everything the composer says about it is not.

\subsection*{What it is suppressing}

Start with the control, because the directive is easier to understand as a
silencer than as an instruction. The same pair with no \code{@override} anywhere:

\begin{minted}[fontsize=\footnotesize,breaklines]{text}
The Object "Product" defines the same fields in multiple subgraphs without the "@shareable" directive:
 The field "availableQuantity" is defined in the following subgraphs: "catalog", "inventory".
 However, it is not declared "@shareable" in any of them.
\end{minted}

Chapter~\ref{ch:composition} spent a section on that error and on how little it
tells you. Add the directive and it goes away. The composer's own walker says
why, in \code{@wundergraph/composition} 0.63.2, and the comment is one line:

\begin{minted}[fontsize=\footnotesize,breaklines]{text}
// overridden fields should not trigger shareable errors
\end{minted}

\index{override@\code{@override}!invisible in the client schema}%
Now the part that surprised me. The composition succeeds and \emph{the
client-facing schema is identical either way}. There is no trace of the
migration in what a client can see, which is correct and is also why nobody
notices a stale \code{@override} for two years. The only place the effect is
visible is the routing table inside the execution config:

\begin{minted}[fontsize=\footnotesize]{text}
without @override    the pair does not compose at all
with    @override    Product.availableQuantity -> inventory
\end{minted}

Catalog still publishes the field in its own schema. The router will never ask
it for one. That is the whole mechanism, and the case asserts it by reading the
datasource configurations rather than by reading the client schema, because the
client schema cannot see it.

\subsection*{Three ways to get it wrong}

\index{override@\code{@override}!from: naming yourself}%
Name your own subgraph in \code{from:} and the composer is direct about it:

\begin{minted}[fontsize=\footnotesize,breaklines]{text}
Cannot override field "Product.availableQuantity" because the source and target subgraph names are both "inventory"
\end{minted}

\index{override@\code{@override}!declared twice}%
Have two subgraphs claim the same field from a third and it is direct again,
although it has clearly been assembled from parts:

\begin{minted}[fontsize=\footnotesize,breaklines]{text}
The "@override" directive must only be declared on one single instance of a field. However, an "@override" directive
was declared on more than one instance of the following field: " The field "Product.availableQuantity" declares an
@override directive in the following subgraphs: "pricing", "inventory".".
\end{minted}

The nested quotes and the stray space are the composer's, not mine: the second
sentence is interpolated already formatted. The case asserts it as written
rather than as intended, because a gate that quietly normalises a message is a
gate that will not notice when the message changes.

\index{override@\code{@override}!a misspelled subgraph name}%
And then the one that costs an afternoon. Misspell the subgraph you are taking
the field from, while the field really is in two subgraphs:

\begin{minted}[fontsize=\footnotesize,breaklines]{text}
The Object "Product" defines the same fields in multiple subgraphs without the "@shareable" directive:
 The field "availableQuantity" is defined in the following subgraphs: "catalog", "inventory".
 However, it is not declared "@shareable" in any of them.
\end{minted}

The override did not apply, so the duplication came back, so the composer
reports the duplication. The word \code{katalog} appears nowhere. Neither does
\code{@override}. You have just added a directive, the build has just started
failing, and the error is about a directive you did not add and a type rather
than a service. The case asserts both absences, because the absence is the
finding.

This is chapter~\ref{ch:composition}'s thesis with a new cause: three different
mistakes there all reported as a shareability error on the key field, and here
is a fourth mistake reporting as a shareability error on an ordinary one.

\subsection*{The book's first composition warning}

\index{composition!warnings!the first one in this book}%
Chapter~\ref{ch:composition} had nothing to say about
\code{--suppress-warnings}, because Mosaic produced no warnings for it to
suppress. Chapter~\ref{ch:entity-resolution-done-right} could only report that
the flag does not touch an error. Here is a warning.

\code{@override(from: "legacy-monolith")} where no subgraph by that name exists,
and where nothing else declares the field, so there is no duplication to report
instead:

\begin{minted}[fontsize=\footnotesize,breaklines]{text}
The following warnings were produced while composing:
The Object type "Product" defines the directive "@override(from: "legacy-monolith")" on the following field:
"availableQuantity".
The required "from" argument of type "String!" should be provided with an existing subgraph name.
However, a subgraph by the name of "legacy-monolith" does not exist.
If this subgraph has been recently deleted, remember to clean up unused "@override" directives that reference this
subgraph.
\end{minted}

Exit code 0. The config is written and the graph is unchanged. Note the last
sentence, which is the composer telling you what this warning is actually for:
the ordinary way to reach it is not a typo but a subgraph that was deleted after
its fields were taken, leaving directives behind that point at nothing.

Compose again with \code{--suppress-warnings} and it prints nothing and writes a
\emph{byte-identical} config. So the flag does exactly what its name says, which
is the first time in three chapters that anything has been able to demonstrate
it. The case asserts the two files match rather than asserting two hashes I
worked out by hand.

\subsection*{Progressive \code{@override}, and where this stack stops}

\index{override@\code{@override}!progressive}%
Apollo Federation 2.7 added a second argument. A label lets you move a
percentage of traffic for one field rather than all of it, and this is Apollo's
own example of one~\autocite{apollo2026migratefields}:

\begin{graphqlsdl}[fontsize=\footnotesize]
description: String @override(from: "BarSubgraph", label: "percent(1)")
\end{graphqlsdl}

HotChocolate has it. \code{OverrideAttribute} takes
\code{(string from, string? label = null)}, and the library's own test builds a
2.7 schema and snapshots the two-argument directive definition. What gates it is
an internal attribute whose comment is the clearest statement of the rule
anywhere:

\begin{minted}[fontsize=\footnotesize,breaklines]{csharp}
// prior to version 2.7 @override only specified "from" parameter
if (descriptor.GetFederationVersion() < FederationVersion.Federation27)
{
    var desc = (IDirectiveTypeDescriptor<OverrideDirective>)descriptor;
    desc.BindArgumentsExplicitly();
    desc.Argument(t => t.From);
}
\end{minted}

\code{FederationVersion.Default} is \code{Federation26}, which is why every
schema in this book emits the one-argument form and why
chapter~\ref{ch:how-a-federated-query-actually-runs} measured v2.6 in the
published SDL.

So opt a subgraph into 2.7 and write a label. The composer has an opinion:

\begin{minted}[fontsize=\footnotesize,breaklines]{text}
The subgraph "inventory" could not be federated for the following reason:
The 1st instance of the directive "@override" declared on coordinates "Product.availableQuantity" is invalid for the
following reason:
 The definition for "@override" does not define the following argument that is provided: "label".
\end{minted}

\index{wgc@\code{wgc}!its own @override definition}%
Exit 1, no config written. The reason is one constant, in
\code{@wundergraph/composition} 0.63.2:

\begin{minted}[fontsize=\scriptsize,breaklines]{javascript}
// directive @override(from: String!) on FIELD_DEFINITION
exports.OVERRIDE_DEFINITION = {
    arguments: [
        {
            kind: graphql_1.Kind.INPUT_VALUE_DEFINITION,
            name: (0, utils_1.stringToNameNode)(string_constants_1.FROM),
            type: {
                kind: graphql_1.Kind.NON_NULL_TYPE,
                type: (0, utils_1.stringToNamedTypeNode)(string_constants_1.STRING_SCALAR),
            },
        },
    ],
\end{minted}

One argument. \textbf{Progressive \code{@override} is unavailable on this
stack}, and
it fails at composition rather than at the router, which is the better of the
two places for it to fail. Apollo documents the feature as a GraphOS Enterprise
one and notes that the router's plan cache keys on the set of unique overridden
labels in an operation~\autocite{apollo2026migratefields}, so a graph with many
labels is a graph with many plans; that is a cost worth knowing about even where
the feature is available.

What is left when it is not available is the all-or-nothing form, and the honest
answer about what that costs is that it moves the risk rather than removing it.
You cannot ramp traffic, so the ramp has to happen somewhere else: behind a
feature flag inside the new resolver, or by running the new subgraph against
production data and comparing before you publish its schema at all. Both are
work the directive was supposed to save you.

\medskip
That is the mechanism for moving one field. What the finished graph does with a
query is a different picture, and the router draws it.

\index{override@\code{@override}|)}%
